\chapter{Both uniform C1 source charts on the entire bounded \(Z=1\) state chart}
\label{component:pa4f6d1e57f34}
\noindent\textbf{Component }\path{pa4f6d1e57f34}\par
{\small\noindent Source: \path{NEW/root/uniform_c1_z_charts_rederived.tex}\par}
{\small\noindent SHA256: \path{a4f6d1e57f340a918497f38bffc46735a5d8d0be05ef6da8c8f6b37cea682186}\par}
{\small\noindent Inventory use: fresh complete EZ1--19 Z-state charts\par}
\medskip
% Fresh whole-parent rederivation, not an unavailable historical witness.
\section{Both uniform C1 source charts on the entire bounded \(Z=1\) state chart}
\label{pa4f6d1e57f34:sec:uniform-c1-z-new}

\paragraph{Domain and conclusion.}
Use the original physical uniform C1 source with $d=1$,
$\lambda=n=0$, $k,f,g,e=kf/g$ positive infinitesimals, and
all original nonnegative event weights. Retain
\[
 r=1-af^2,\quad R=1+k^2r,\quad
 T=1+k^2-tk^2g^2,\quad M=R-t>0,\quad J=T-a>0,
 \qquad a,t>0.                                                   \tag{EZ1}
\]
In particular $a,t$ are bounded, $r,R,T$ are positive units,
and $M,J$ may approach zero at arbitrary rates.
For every independent testing $u,v,X,Y$ infinitesimal and
$Z=1$, $W$ arbitrary bounded of either sign, the entire original
parent is a separated tensor SOS. This includes both bounded
and unbounded $b=g/k$, and every permitted $\beta=0$ face.
It does not claim the complementary $W=1,Z$-infinitesimal chart.

\subsection{The exact source and a common set of bounded coefficients}

Use the seven short native coquotients
$c_2,c_\beta,c_3,c_5,c_4,c_\alpha,c_d$ in the retained
\path{endpoint_source.tex}, now with $\eta=f^2,\zeta=g^2$.
Their residues are $2,2,2,2,2,2,4$.
In the present $d=1$ normalization the exact weights are
\[
\begin{gathered}
 C=\frac{c_3}{af^2c_d},\quad A=\frac{f^2c_4}{c_d},\quad
 m=\frac{k^2c_2}{g^2Mc_d},\quad
 \alpha=\frac{k^4g^2c_\alpha}{c_d},\\
 D+\tfrac12=\frac{c_5}{Jf^2c_d},\qquad
 \beta+\tfrac12=\frac{k^2c_\beta}{g^2tc_d}.
\end{gathered}                                                   \tag{EZ2}
\]
Every use below is of these actual weights. In particular
$\beta\geq0$ is retained, not replaced by a formal source condition.

The following two pure-source coordinate choices are exhaustive:
\[
\begin{array}{c|ccccc}
 &s&\omega&P&G&\kappa=P/G\\ \hline
 b\text{ bounded}&e&b^2&k^2&k^2&1\\
 b\text{ unbounded}&f&1&k^2&g^2&(k/g)^2
\end{array}                                                     \tag{EZ3}
\]
In both rows $f^2=s^2\omega$, $g^2=\omega G$,
$s,\sqrt P,\sqrt G$ are infinitesimal, and $\omega,\kappa$
are positive bounded. In the second row $\kappa$ is infinitesimal.
Here $\beta\geq0$ gives exactly
$t\leq2\kappa c_\beta/c_d$, so $t=\kappa\tau$ with
$\tau>0$ bounded; consequently $M$ is a unit in this row.

Define bounded source coefficients
\[
 c_*=\frac{c_3}{Tc_d},\quad d_*=\frac{c_5}{Tc_d},\quad
 a_*=\frac{c_4}{c_d},\quad \gamma_*=\frac{c_\alpha}{c_d},
 \qquad B_*=\omega\beta t
              =\frac{\kappa c_\beta}{c_d}-\frac{\omega t}{2}.
                                                               \tag{EZ4}
\]
The first four have residue $1/2$; $B_*\geq0$ may vanish.
Use the further table
\[
\begin{array}{c|ccc}
 &\mu&\chi&\psi\\ \hline
 b\text{ bounded}&c_2/(Rc_d)&t&tM\\
 b\text{ unbounded}&\kappa c_2/(Mc_d)&0&0 .
\end{array}                                                     \tag{EZ5}
\]
All entries are bounded. In the first row $\mu$ has residue
$1/2$; in the second $\mu$ is infinitesimal.

\subsection{Exact removal of the source poles}

The original quarter-parent is
\[
\begin{split}
 F={}&mr_2^2+CD_vr_3^2+AD_uD_vY^2
 +(D+\tfrac12)r_5^2+(\alpha+\tfrac12)D_u(vY-1)^2\\
 &+(\beta+\tfrac12)D_v(X+u)^2-\tfrac12D_uD_v,
\end{split}                                                     \tag{EZ6}
\]
where $r_2=vW-X+uvY-u$, $r_3=uW-Y$,
$r_5=uW+uvX-Y-v$, $D_u=1+u^2$, $D_v=1+v^2$.
Put
\[
 z=sW,\qquad q_0=f^2uW^2=\omega uz^2,\qquad
 z_0=f^2uW=s\omega uz.
\]
Retain the following complete nonnegative packets of the original
polynomial:
\[
 \frac{CJ}{T}D_v(r_3-az_0)^2,\qquad
 \frac{(D+1/2)a}{T}(r_5+Jz_0)^2,\qquad
 \beta D_v(X+u-tq_0)^2.                                        \tag{EZ7}
\]
For bounded $b$, also retain
\[
                       \frac{mt}{R}(r_2-Mq_0)^2.              \tag{EZ8}
\]
Each vanishes at the actual source contact. They are literal
positive squares, even when $\beta=0$ or a source opening is small.
The first two removals use $a+J=T$ and the exact identity
\[
 \frac aT[(L+Jq)^2-TJq^2]+\frac JT(L-aq)^2=L^2.
\]
The last removal uses
$r_2^2-(t/R)(r_2-Mq_0)^2
=(M/R)(r_2^2+2tq_0r_2-tMq_0^2)$.
There is no mixed denominator.

Use all three unrestricted normal coordinates
\[
 N_x=X+u-tq_0,\qquad
 N_y=(Y-ruW)/s,\qquad N_v=(v-f^2uW)/s.
                                                               \tag{EZ9}
\]
Only $N_x$ is necessarily bounded. The exact bounds needed later are
$N_x,sN_y,sN_v\in\mathrm I$; they follow from the stated testing
chart and bounded $W$. No boundedness of $N_y,N_v$ is assumed.
Define
\[
 X_*=-u+tq_0+N_x,\qquad
 y_*=ruz+s^2N_y=sY,\qquad
 v_*=\omega uz+N_v=v/s,
\]
\[
 L_2=zv_*+uv_*y_*-X_*-u,\quad
 L_3=T\omega uz-N_y,\quad
 L_5=-N_y+v_*(uX_*-1).
\]
After multiplying the original quarter-parent by the positive
pure-source $\omega$ and removing precisely the scaled packets
EZ7--EZ8, the \emph{entire} remaining polynomial is
\[
\begin{split}
\mathcal F={}&
 \mu(L_2^2+2\chi q_0L_2-\psi q_0^2)
 +c_*D_v(L_3^2-TJ\omega^2u^2z^2)\\
 &+\omega^2a_*D_uD_vy_*^2
 +d_*(L_5^2-aT\omega^2u^2z^2)\\
 &+(\omega^2P^2G\gamma_*+\omega/2)
                         D_u(v_*y_*-1)^2
 +\frac{\omega}{2}D_v(X_*+u)^2\\
 &+B_*D_v\{2q_0(X_*+u)-tq_0^2\}
                         -\frac{\omega}{2}D_uD_v,
 \qquad D_v=1+s^2v_*^2 .
\end{split}                                                     \tag{EZ10}
\]
This is a short full-parent definition of every scalar and normal
coefficient used below. In particular no normal was set to zero
in obtaining it. All its coefficient arrays are bounded through
$a,J,M,t$; the unbounded original weights occur only in the
retained complete squares.

\subsection{The entire scalar is a cubic circuit with finite corrections}

Set $U=u^2,V=z^2$. The polynomial obtained from EZ10 at
$N_x=N_y=N_v=0$ is exactly
\[
 \mathcal F_0=\omega^2\Phi(U,V),\qquad
 \Phi=\gamma_*P^2G(1+U)+UV\,H(U,V),                            \tag{EZ11}
\]
where $H$ has bounded finite polynomial coefficient arrays.
The apparent divisions by $\omega^2$ and $UV$ in this definition
are coefficientwise polynomial extractions: substitution in EZ10
shows those factors before any evaluation. No testing coordinate
is inverted.

At the actual source contact $u^2=P,z^2=G$ all three normal
coordinates and all retained packets vanish. The original contact
value and stationary equations give
\[
 \Phi(P,G)=\Phi_U(P,G)=\Phi_V(P,G)=0.                            \tag{EZ12}
\]
The exact native coquotients in EZ2 independently verify these
identities. Substituting EZ11, and cancelling nonzero pure-source
factors in the source field, gives
\[
 H(P,G)=-\gamma_*P(1+P),\quad
 H_U(P,G)=\gamma_*,\quad
 H_V(P,G)=\gamma_*\kappa(1+P).
\]
Write $\Delta=U-P$, $\Gamma=V-G$.
Finite monic division, first in $U-P$ and then in $V-G$,
therefore gives the complete scalar identity
\[
\begin{split}
 \Phi={}&\gamma_*\{\mathcal M+P\kappa U\Gamma^2\}\\
 &+UV\{\Delta^2H_{20}+\Delta\Gamma H_{11}
                                      +\Gamma^2H_{02}\},\\
 \mathcal M={}&U^2V+\kappa UV^2+P^2G-3PUV .
\end{split}                                                     \tag{EZ13}
\]
All quotients are finite bounded coefficient polynomials; no
inverse of $P,G$ is introduced by their monic construction.
The extra term $\gamma_*P\kappa U\Gamma^2$ is retained as a
nonnegative square.

In the inverse-$b$ row one further coefficient factor is needed:
\[
                         H_{02}=\kappa\widetilde H_{02},
 \qquad\widetilde H_{02}\text{ has bounded arrays}.              \tag{EZ14}
\]
Here is a finite source-division certificate, rather than an
estimate involving independent testing rates. Write
$P=\kappa G$, $t=\kappa\tau$ before substituting the native
coquotients into EZ10. Every denominator is among the pure-source
units $M,T,c_d$ and their factors. At formal $\kappa=0$,
\[
 P=t=\mu=B_*=0,\quad
 c_*=d_*=\frac{1+(1-a)s^2}{2},\quad
 a_*=\frac1{2(1-as^2)} .
\]
The formula gives $H(0,V)=0$ identically in $V$.
After clearing unit denominators, ordinary polynomial division
by $\kappa$ therefore has zero remainder. Apply the two monic
$V-G$ divisions to this quotient. This proves EZ14 with bounded
arrays. In the bounded-$b$ row $\kappa=1$, so it needs no such
factor.

\subsection{Six explicit strict roots and their conductors}

The following identities apply to both rows of EZ3. Put
\[
 \mathsf h=G^{1/6},\quad\mathsf c=\kappa^{1/6},\quad
 \mathsf r=\sqrt[3]u,\quad\mathsf s=\sqrt[3]z,
\]
\[
 (L_1,L_2,L_3)=
 (\mathsf r^4\mathsf s^2,\
   \mathsf c^2\mathsf r^2\mathsf s^4,\
   \mathsf h^6\mathsf c^4),\qquad
 (f_1,f_2,f_3)=
 (\mathsf r^2\mathsf s,\
   \mathsf c\mathsf r\mathsf s^2,\
   \mathsf h^3\mathsf c^2).
\]
Define $R_{2i-1}=f_i(L_1-L_2)$ and
$R_{2i}=f_i(L_1-L_3)$ for $i=1,2,3$.
Then, exactly,
\[
 \mathcal M=\sum_{i=1}^3
       (R_{2i-1}^2-R_{2i-1}R_{2i}+R_{2i}^2).                  \tag{EZ15}
\]
Each pair has the fixed strictly positive Gram
$\left(\begin{smallmatrix}1&-1/2\\-1/2&1\end{smallmatrix}\right)$.
Real cube roots of signed testing quantities are permitted:
$z=sW$ has the separated cube root
$s^{1/3}\sqrt[3]W$. The $L_i=f_i^2$ identity is valid for
either sign of $u,z$. No mixed radical occurs.

For explicit finite allocation define
\[
 J_1=\mathsf r^3R_1+\mathsf c^2\mathsf r\mathsf s^2R_2,\qquad
 J_2=-\mathsf h^6R_1+\mathsf s^6R_2
                              -\mathsf h^3\mathsf s^3R_5.
\]
Direct multiplication gives
\[
\begin{aligned}
 uz\Delta&=J_1,&u^2z\Gamma&=J_2,\\
 \kappa uz\Gamma
   &=\mathsf c^2\mathsf r\mathsf s^2(R_2-R_1)
                                  -\mathsf c^3\mathsf s^3R_3,\\
 UV\Delta^2&=J_1^2,\\
 UV\Delta\Gamma
   &=\mathsf s^6R_2^2-\mathsf h^6R_1^2
             -\mathsf h^3\mathsf s^3R_1(R_5+R_6),\\
 \kappa UV\Gamma^2
   &=\mathsf s^6(2R_1^2-3R_1R_2+R_2^2+R_3^2)
                      +\mathsf h^3\mathsf s^3R_1(R_5-R_6).
\end{aligned}                                                   \tag{EZ16}
\]
All coefficients on the right are infinitesimal: $\mathsf r,
\mathsf s,\mathsf h$ are infinitesimal and $\mathsf c$ is
bounded. It is harmless that $\mathsf c=1$ in the first row.
Equations EZ13--EZ16 therefore put the entire scalar, apart
from its retained positive square, on the six rows $\omega R_i$
with a bounded Gram whose real residue is $\tfrac12$ times the
strict matrix in EZ15.

\subsection{Every complete first normal and all higher terms}

Define $L_x,L_y,L_v$ as half the three derivatives of the
whole polynomial EZ10 at zero normals. Expanding its short
factored form gives exactly
\[
 L_x=\omega uz^2T_x(U,V),\qquad
 L_y=\omega uzT_y(U,V),\qquad
 L_v=\omega uzT_v(U,V),                                       \tag{EZ17}
\]
where each $T_i$ is a polynomial with bounded coefficient arrays.
This supplies a finite construction of every coefficient:
differentiate EZ10, set the three normals to zero, and extract
the displayed monomial. No derivative is dropped.
The exact actual-source contact equations give $T_i(P,G)=0$.
Consequently the monic divisions are
\[
 T_i=(U-P)A_i(U,V)+(V-G)B_i(V).                               \tag{EZ18}
\]
For inverse $b$, each $B_i=\kappa\widetilde B_i$ with
bounded arrays. Indeed the same source substitution used in
EZ14 gives $T_i(0,V)=0$ identically at $\kappa=0$.
Divide the cleared source polynomial by $\kappa$ and then by
the monic $V-G$. This is an exact coefficient construction;
it uses neither $1/u$, $1/z$, nor an evaluated mixed quotient.
For bounded $b$, $\kappa=1$ again suffices.

Combining EZ16--EZ18 proves that both complete state normals
and the complete testing-angle normal have infinitesimal
coefficients on the strict six-row list $\omega R_i$.
The additional factor $z$ in $L_x$ is bounded and infinitesimal.
This argument uses the full first-normal polynomial, not just
its value or linear approximation at contact.

The full normal expansion of EZ10 has a bounded quadratic Gram
whose real residue, in order $(N_x,N_y,N_v)$, is
\[
 \begin{pmatrix}
  \bar\mu+\bar\omega/2&0&0\\
 0&1&1/2\\
 0&1/2&1/2
 \end{pmatrix}\succ0.                                         \tag{EZ19}
\]
For bounded $b$, $\bar\mu=1/2$; for inverse $b$,
$\omega=1$ and $\bar\mu=0$. Both give a strictly positive
first pivot. The other block has pivots $1,1/4$.

Here is also a finite justification for all higher terms, since
the scaled normals $N_y,N_v$ may be unbounded. Expanding EZ10
in the three normals gives degree at most four. A monomial
of normal degree at least three containing $N_x$ has total
$N_y,N_v$ degree at most two. Assign those factors, and any
needed extra $N_x$ factor, to two Gram rows; the remaining
factor contains $N_x\in\mathrm I$. For a higher monomial
without $N_x$, each of the $N_y,N_v$ exponents is at most two,
and its coefficient has a factor $s^2$. If its degree is
three, the unassigned coefficient is a bounded multiple of
$s(sN_y)$ or $s(sN_v)$. If its degree is four, it is a
bounded multiple of $(sN_y)(sN_v)$ (or of one of their
squares). All are infinitesimal by EZ9. These support and
divisibility statements follow directly from the factored
quadratics in EZ10 and are checked on the entire polynomial.
Thus no unbounded scaled normal is hidden in a Gram coefficient.

The scalar, all first normals, and the complete hard block now
give a finite bounded Gram with strictly positive real residue.
The separated strict-Gram lemma supplies tensor SOS.
Restore the whole retained packets EZ7--EZ8 and the positive
square in EZ13, and divide by the pure-source positive $\omega$.
This proves the stated original-parent theorem on both source
rows, without a comparison of source and testing rates.

\paragraph{State and boundary scope.}
The signs of $W$ and of the testing angles were not restricted
in EZ10--EZ19. The polynomial and cube-root identities include
literal zero testing coordinates. If a bounded normalized
$Z=1$ state has $X$ or $Y$ a unit, the original pure-testing
core has positive residue at the double-small angle corner;
its positive source multiple is already a separated square,
and all original event packets are nonnegative squares.
Thus this proof supplies the entire $Z$-unit state chart.
Complex states split into their real and imaginary quadratic
evaluations. The two strong-current outer testing identities
remain separate retained inputs to a future uniform all-tests
join, as does the complementary $W$-unit chart.

\paragraph{Verification boundary.}
The accompanying checker verifies the full packet restorations,
all source contacts and first jets in both native substitutions,
the full first-normal factorizations, the finite source-factor
divisions, the exact cubic conductors, and the whole higher-normal
support and hard residues. The proof of boundedness, strictness
and the stated source/state scope is the analytical argument
above. This result does not by itself establish the global
finite semidefinite-lift theorem.
